% ============================================================
% Chapter 5: Development
%   Section: The News Map, an App for Readers
% Included from chapters/05-development/development.tex
% Source: docs/decisions/news-map.md (built 2026-10-29 on
% feature/news-map, on top of the plan, at the author's request).
% ============================================================

\section{The News Map, an App for Readers}
\label{sec:dev-news-map}

The reader view shows one checked article at a time. On 29~October~2026
the author asked for a second way in, for readers: the news as one
enormous graph, in which every published article is a point joined to
the people, places and topics it names, and the topics sit in sections,
so the whole reads as a cloud. Around it, a search, a page for the
day's news, a selection built from a reader's interests, and accounts
to keep them. It reads only what the reader view already shows anyone:
published articles, their titles, outlets, topics, names and verdicts,
never a body or a scraped page.

\paragraph{A place that never moves.} The map is a disc. The five topic
groups are sectors of it, and each of the 23 topics is an anchor in its
group's sector, on one of two rings. An article sits at its main
topic's anchor, is pulled a short fixed distance towards a second topic
only when that topic's confidence ties the first, and is moved within
its topic's cluster by its own embedding, projected on a fixed random
basis. Its position therefore depends on the article and a layout
version alone: a newly published article moves no other one, which a
test holds. Three measurements shaped this. A confidence-weighted mean
of the anchors, the first design, put nearly every article near the
centre, because the classifier's confidences are compressed (the first
article read had three topics within 0.023 of each other, in three
groups). A pull expressed as a share of the distance carried that
article past a neighbouring anchor. And the first way of squashing the
projection into a cluster made rings and crescents when drawn at 3,000
articles; fitted to the projections of the 139 hand-labelled claims, a
hyperbolic tangent keeps clusters as clouds.

\paragraph{Drawing it.} The points are drawn on the graphics card in a
single call, whatever their number. The library was chosen by
measurement rather than by name: the same synthetic points were drawn
by regl and by sigma.js, a graph library, along the same pan and zoom
in a headless browser, on the laptop's integrated GPU and in software
(Table~\ref{tab:news-map-renderers}). On the GPU both were limited by
the measurement's own read-back; without a GPU, regl was about sixteen
times faster, and sigma.js could not draw 50,000 points.

\begin{table}[h]
  \centering
  \caption{Milliseconds per frame, median / 95th percentile, of the same
  synthetic points drawn by each library (29~October~2026, the
  machine shared with a model benchmark).}
  \label{tab:news-map-renderers}
  \begin{tabular}{lrr}
    \hline
    & regl & sigma.js \\
    \hline
    Integrated GPU, 10,000 points & 8.8 / 14.1 & 14.2 / 23.0 \\
    Integrated GPU, 50,000 points & 10.4 / 19.9 & 15.1 / 25.7 \\
    Software, 10,000 points & 19.9 / 36.6 & 324.6 / 543.6 \\
    Software, 50,000 points & 33.5 / 61.9 & did not finish \\
    \hline
  \end{tabular}
\end{table}

\paragraph{Search, the day, and what a reader follows.} A search finds
names (folded for case, accents and punctuation, then by character
trigrams), words in titles and topic keywords in English and Spanish,
and, when asked, meaning: the query is embedded and compared with the
vectors of published articles only. That last part costs a call to the
inference service, which the fact-checking pipeline shares, so it is
limited per visitor and cached. The day's page brightens one day's
articles on the map, the day being Madrid's. A reader's interests,
saved stories and the connections they draw between stories are kept
in their own database, never in the lake or the graph, which a sync
rebuilds; a reader's connection is drawn on their map alone.

\paragraph{For you, explained.} Each published article is scored in
four parts --- a followed topic, followed names, nearness in meaning to
the saved stories, and freshness --- with weights that sum to one, so
the parts shown beside a story add up to its score. The weights are
reasoned, not fitted: there is no labelled set of what readers wanted.
Within the first ten, at most seven come from a topic the reader
follows and three from any other, and ``not interested'' lowers, for
that reader alone, the part that counted most.

\paragraph{Accounts.} Passwords are kept as argon2id hashes and session
tokens as their SHA-256, so a copy of the database opens nothing.
Signing in is limited per name and per client and takes the same work
whether or not the name exists. Every change a signed-in reader makes
carries a token only their own page can read, and a reader can export
or delete everything kept about them at once.

\paragraph{Its budget.} Measured in a headless browser on the laptop's
integrated GPU with synthetic maps placed by the real layout, a map of
10,000 stories is drawn from a warm cache in 373~ms, reaches the browser
as 354~KB, and finds the point under the pointer in 0.1~ms; 50,000
stories take 499~ms and 1.77~MB. The first measurement found the
payload uncompressed (2.09~MB at 10,000), because the frontend's proxy
undid the backend's compression; it is compressed again there now.
\pending{Whether panning and zooming hold 60 frames a second: measured
under load on 29~October, dropped frames fell as the map grew (7.2\% at
2,000 stories, 0.8\% at 50,000), which points at the machine rather
than the drawing. The run on an idle machine goes here
(\texttt{scripts/bench\_news\_map.py}).}
